% ===========================================================================
%  Appendice A — Configurazioni complete
% ===========================================================================
\chapter{Configurazioni complete}
\label{app:configurazioni}

Questa appendice riporta in forma tabellare il contenuto dei file
\file{config.json} salvati dalle run utilizzate nella tesi. Ogni
\file{config.json} è autodescrittivo: include sia i campi della dataclass sia le
proprietà calcolate (\cref{sec:base-config}), e può quindi essere interpretato
senza reistanziare il codice.

\section{Esperimento 01 --- preset di \texorpdfstring{\classe{GPTConfig}}{GPTConfig}}

\begin{table}[htbp]
\centering
\caption[Preset di configurazione dell'esperimento 01]{%
Preset disponibili in \file{01\_quantum\_gpt/src/config/}. Il valore di $\nq$ è
calcolato come $\demb/\nhead$ e non è un campo indipendente. Il preset
\code{big} non è utilizzabile nella variante quantistica: 64 qubit non sono
simulabili.}
\label{tab:preset-gpt}
\small
\begin{tabular}{@{}l r r r r r r r c@{}}
\toprule
\textbf{Preset} & $\demb$ & $\nhead$ & $n_{\text{layer}}$ & $\nq$ & \code{block\_size} & \code{max\_iters} & \code{lr} & \code{use\_quantum} \\
\midrule
\code{default} & 32  & 8  & 4 & 4  & 64  & \num{5000}  & \num{1e-3} & \code{False} \\
\code{fast}    & 8   & 2  & 2 & 4  & 8   & 500         & \num{3e-3} & \code{True} \\
\code{thesis}  & 8   & 2  & 2 & 4  & 16  & \num{1500}  & \num{3e-3} & \code{False} \\
\code{light}   & 32  & 4  & 2 & 8  & 16  & \num{1000}  & \num{1e-3} & \code{False} \\
\code{heavy}   & 128 & 32 & 9 & 4  & 128 & \num{20000} & \num{3e-4} & \code{False} \\
\code{big}     & 384 & 6  & 6 & 64 & 256 & \num{20000} & \num{3e-4} & \code{False} \\
\bottomrule
\end{tabular}
\end{table}

La campagna definitiva usa il preset \code{thesis}, riportato per esteso qui
sotto; le due varianti differiscono soltanto nel campo \code{use\_quantum}.

\begin{center}
\begin{tabular}{@{}l l l l@{}}
\toprule
\textbf{Campo} & \textbf{Valore} & \textbf{Campo} & \textbf{Valore} \\
\midrule
\code{tokenizer\_class} & \code{CharTokenizer} & \code{batch\_size} & 8 \\
\code{block\_size} & 16 & \code{max\_iters} & 1500 \\
\code{eval\_interval} & 100 & \code{eval\_iters} & 100 \\
\code{learning\_rate} & \num{3e-3} & \code{dropout} & 0 \\
\code{n\_embd} & 8 & \code{n\_head} & 2 \\
\code{n\_layer} & 2 & \code{n\_qubits} & 4 \\
\code{n\_qlayers} & 2 & \code{q\_device} & \code{default.qubit} \\
\code{device} & \code{cpu} & \code{use\_quantum} & \code{False}/\code{True} \\
\bottomrule
\end{tabular}
\end{center}

\section{Esperimento 02 --- \texorpdfstring{\classe{QViTConfig}}{QViTConfig}}

\begin{center}\small
\begin{tabular}{@{}l l l l@{}}\toprule
\textbf{Campo}&\textbf{Valore}&\textbf{Campo}&\textbf{Valore}\\\midrule
\code{dataset}&\code{pathmnist}&\code{max\_epochs}&50\\
\code{embed\_dim}&8&\code{n\_head}&2\\
\code{n\_layer}&2&\code{ffn\_dim}&32\\
\code{patch\_size}&7&\code{batch\_size}&128\\
\code{dropout}&0,1&\code{weight\_decay}&\num{1e-4}\\
\code{n\_qlayers}&2&\code{q\_device}&\code{default.qubit}\\
\code{device}&\code{cpu}&\code{use\_quantum}&\code{False}/\code{True}\\\bottomrule
\end{tabular}\end{center}

\section{Esperimento 03 --- \texorpdfstring{\classe{AblationConfig}}{AblationConfig}}

\begin{center}\small
\begin{tabular}{@{}l l l l@{}}\toprule
\textbf{Campo}&\textbf{Valore}&\textbf{Campo}&\textbf{Valore}\\\midrule
\code{datasets}&Path/Blood/Derma&\code{train\_subset}&100\\
\code{optimizer\_steps}&400&\code{eval\_interval}&40\\
\code{embed\_dim}&10&\code{n\_head}&2\\
\code{n\_layer}&2&\code{ffn\_dim}&40\\
\code{patch\_size}&7&\code{batch\_size}&128\\
\code{dropout}&0&\code{weight\_decay}&0\\
\code{n\_qlayers}&2&\code{q\_device}&\code{default.qubit}\\
\code{device}&\code{cuda}&\code{evaluate\_noise}&\code{False}\\\bottomrule
\end{tabular}\end{center}
Le varianti sono \code{vanilla}, \code{bounded\_mlp} e
\code{quantum\_reg}; canali e classi sono determinati dal dataset.

\section{Esperimento 04 --- argomenti della riga di comando}

L'esperimento~04 non usa una dataclass di configurazione: i suoi parametri sono
argomenti della CLI, registrati nel campo \code{command} di
\file{run\_manifest.json} e replicati nei campi di primo livello di
\file{results.json} (\code{dataset}, \code{seed}, \code{n\_qubits},
\code{n\_class0}, \code{n\_class1}, \code{q\_device}).
